\paragraph{C.3aa (Refuted): every non-top positive-codimension local finiteness assertion over $\mathbb C$.}
For every pair of integers
\[
                        N\ge4,\qquad 2\le r\le N-2,
\]
there is a specified excellent normal strictly henselian local
$\mathbb C$-algebra $R$ of dimension $N$, with regular puncture
$X$, such that
\begin{equation}\label{eq:C3aa-nonvanishing}
                         B^r(X)_{\mathbb Q}\ne0.
\end{equation}
Here $B^r$ is cycles modulo ordinary algebraic equivalence,
and codimension is inherited from the local spectrum as in
C.3e. In particular the integral group $B^r(X)$ is not
finite. Together with the divisor examples C.3v and C.3y,
this settles every non-top positive-codimension pair
$1\le r\le N-2$ over $\mathbb C$. This is a nonvanishing
assertion; no infinite-dimensionality claim for these
higher rational groups is made. The remaining top pair
$r=N-1$, $N\ge3$, is not decided by this argument.

\begin{proof}
\emph{1. Fix the explicit Nori witness in the smallest needed dimension.}
Apply the fully proved construction C.1i with its indices
$c=r$ and $n=r+1$. Its strict inequalities hold because
$r\ge2$. More explicitly, take the split quadric
\[
 Q=V_+\left(\sum_{i=0}^{r}x_i y_i\right)
           \subset\mathbb P^{2r+1}_{\mathbb Q},
 \qquad H_Q=c_1(\mathcal O_Q(1)),
\]
with the two specified maximal planes
\[
 \Pi_+=V_+(x_0,\ldots,x_r),\qquad
 \Pi_-=V_+(y_0,x_1,\ldots,x_r).
\]
Set $d=4r+2$. Let $T\subset Q$ be the geometric generic
complete intersection of $r-1$ degree-$d$ sections,
with exactly the independent coefficients and the one
normalized coefficient specified in C.1i. Thus the
defining field is the algebraic closure of that explicit
purely transcendental extension of $\mathbb Q$.
Fix any embedding of this field into $\mathbb C$;
C.1i proves persistence under every such embedding.
Writing $j:T\hookrightarrow Q_{\mathbb C}$ and
$H=j^*H_Q$, the actual codimension-$r$ cycle
\[
              z=j^*([\Pi_+]-[\Pi_-])
\]
satisfies
\begin{equation}\label{eq:C3aa-nori-witness}
    \dim T=r+1,\qquad
    z\notin H B^{r-1}(T)_{\mathbb Q}.
\end{equation}
In particular $z\ne0$ in $B^r(T)_{\mathbb Q}$.
This invokes the exact Nori--Nagel connectivity theorem
and the complete trace/incidence proof already checked
in C.1i; it introduces no new unspecified degree bound.

\emph{2. Raise the local dimension while retaining the obstruction.}
Put
\[
 m=N-r-2\ge0,\qquad V=\mathbb P^m_{\mathbb C}\times T,
 \qquad
 L=\mathcal O_{\mathbb P^m}(1)\boxtimes\mathcal O_T(1).
\]
Then $\dim V=N-1$. If $m=0$, the first factor is a point
and $L=\mathcal O_T(1)$, as intended. Denote the projective
hyperplane class by $h$ and continue to write $H$ for
its pullback from $T$; thus $c_1(L)=h+H$.

The projective-bundle formula, and preservation of
ordinary algebraic equivalence by its pullbacks,
pushforwards and Chern-class operations, give the
graded-ring identity
\begin{equation}\label{eq:C3aa-projective-ring}
 B^*(V)_{\mathbb Q}
        =B^*(T)_{\mathbb Q}[h]/(h^{m+1}).
\end{equation}
For completeness, the inverse to the expression
$\sum_{a=0}^m h^a\alpha_a$ is given by pushing
$h^{m-a}$ times a class to $T$. These operations
preserve ordinary algebraic equivalence, so the
usual Chow decomposition descends to this quotient
without imposing a stronger equivalence relation.
Imposing $h+H=0$ in \eqref{eq:C3aa-projective-ring}
therefore gives exactly
\begin{equation}\label{eq:C3aa-segre-quotient}
 \frac{B^r(V)_{\mathbb Q}}
      {(h+H)B^{r-1}(V)_{\mathbb Q}}
 \simeq
 \frac{B^r(T)_{\mathbb Q}}
      {H^{m+1}B^{r-m-1}(T)_{\mathbb Q}}.
\end{equation}
The map is evaluation at $h=-H$; explicitly it sends
$\sum_a h^a\alpha_a$ to $\sum_a(-H)^a\alpha_a$.
If $r-m-1<0$, the denominator on the right is zero.
Otherwise
\[
 H^{m+1}B^{r-m-1}(T)_{\mathbb Q}
           \subset H B^{r-1}(T)_{\mathbb Q}.
\]
Consequently the pullback of $z$ has nonzero image
in \eqref{eq:C3aa-segre-quotient}, by
\eqref{eq:C3aa-nori-witness}. This argument covers
all $m\ge0$, including $m\ge r$, where the additional
denominator has disappeared entirely.

\emph{3. Verify the cone hypotheses and pass to its strict puncture.}
The hyperplane embedding of $T$ is a smooth complete
intersection in projective space. Its homogeneous
coordinate ring is Cohen--Macaulay, its affine cone
is regular away from the vertex, and the vertex has
codimension $r+2\ge4$. Serre's criterion therefore
makes that ring normal. The complete-intersection
Koszul resolution, equivalently the usual vanishing
of intermediate projective-space cohomology, identifies
it with the full section ring
\[
                    \bigoplus_{n\ge0}H^0(T,\mathcal O_T(n)).
\]
It is generated in degree one. The full section ring
of $(V,L)$ is the Segre product
\[
 A=\bigoplus_{n\ge0}
       H^0(\mathbb P^m,\mathcal O(n))
          \otimes_{\mathbb C} H^0(T,\mathcal O_T(n)).
\]
It is generated in degree one because both factors
are: products of paired degree-one generators span
each displayed tensor product. It is normal, either
by the full-section-ring valuation criterion on the
smooth normal variety $V$, or directly as follows.
The tensor product of the two normal coordinate
rings is the coordinate ring of their product,
which is normal over the perfect algebraically
closed field. The Segre ring is its weight-zero
subring for the action with weights $+1$ and $-1$.
A weight-zero fraction integral over this subring
is integral over the normal tensor ring and hence
belongs to it; being fixed by the action, it lies
in the weight-zero subring. This proves normality.
Thus $L$ is a very ample projectively normal
degree-one-generated embedding of precisely the
kind required in C.3e.

Let $v$ be the vertex of $\operatorname{Spec}A$,
and define the actual local ring and puncture by
\[
 R=\mathcal O_{\operatorname{Spec}A,v}^{\mathrm{sh}},
 \qquad X=\operatorname{Spec}R\setminus\{\mathfrak m\}.
\]
Then $R$ is excellent, normal and strictly henselian,
$\dim R=N$, and $X$ is regular: the cone off its
vertex is the $\mathbb G_m$-bundle of $L^{-1}$
over the smooth variety $V$, and regularity persists
under strict henselization. The integral strict-cone
injection of C.3e, tensored with the flat module
$\mathbb Q$, embeds the left-hand side of
\eqref{eq:C3aa-segre-quotient} into $B^r(X)_{\mathbb Q}$.
The specified pullback of $z$ remains nonzero,
which proves \eqref{eq:C3aa-nonvanishing} and gives
the asserted actual counterexample for each $(N,r)$.
\end{proof}

The algebraic elimination in
\eqref{eq:C3aa-segre-quotient} is also recorded by
\texttt{code/verify\_segre\_nori\_quotient.py}. It checks
the exact identity
\[
 h^a-(-H)^a=(h+H)\sum_{b=0}^{a-1}h^{a-1-b}(-H)^b
\]
for its stated finite sample and verifies the index
and degree inequalities over a finite grid of $(N,r)$.
The universal quotient formula is proved above;
the finite program is an independent sign/index
check and does not compute any Griffiths group.
